\section{Fresh Tuples}
\label{sec:fresh}

This section explains the proof of \cref{thm:fresh}: a minimal fresh tuple
of a table with $k$ starts and $N$ states has bounded total length. We fix
such a table and a fresh tuple $W=(w_1,\dots,w_k)$. Full proofs are in
\cref{app:fresh}.

\subsection{The Idea}\label{sec:fresh-idea}

To shorten $W$ we delete a factor between two cuts of one word that carry
the same state. This preserves endpoints and visited states, but in general
not the channels. A channel compares the rest of a target word with a factor
of a host word. If we delete a factor from the host in the middle of such a
comparison, the two sides no longer match; and after the deletion the rest
of some word may match a factor of the edited host that it did not match
before, at cuts whose states were never compared. Both effects disappear if
the deletion is \emph{invisible to every comparison}: whenever a comparison
runs over a deleted factor on one side, the same factor is deleted at the
corresponding place on the other side.

This suggests deleting a repeated factor at \emph{all} places that look
alike, in all words at once. The question is what ``look alike'' should
mean. Comparing bounded windows of letters is not enough, because a
comparison can run over arbitrarily long equal texts. We therefore cut all
words into blocks by a rule that depends only on the text, so that equal
texts are cut in the same way, up to a bounded number of blocks at their left
ends. The cutting rule is given by \emph{markers}, suffixes of the words
$w_i$ themselves, and it is iterated: the blocks of one level are grouped
into the blocks of the next level, until every word is a single block.
Deleting equal stretches of blocks at equal distances from the right ends of
equal block sequences is then invisible to comparisons, and a pigeonhole
argument finds such stretches whenever some block consists of too many
blocks of the level below.

\subsection{The Marker Hierarchy}\label{sec:fresh-hierarchy}

\begin{definition}[Markers and anchors]\label{def:markers}
For each \emph{level} $t\ge1$ we define a set $\Act_t$ of \emph{active}
coordinates and, for $l\in\Act_t$, a \emph{marker} $\mk tl$, a nonempty
suffix of $w_l$.
\begin{itemize}
\item $\Act_1$ consists of the coordinates with a nonempty word, and $\mk 1l$
  is the last letter of $w_l$.
\item A \emph{level-$t$ anchor} of a text $u$ is a cut $e\ge1$ of $u$ at
  which $\mk tm$ ends, for some $m\in\Act_t$. The level-$t$ anchors of the
  word $w_l$ are its anchors as a text together with its end $\len{w_l}$.
\item A coordinate $l\in\Act_t$ belongs to $\Act_{t+1}$ if $w_l$ has a
  level-$t$ anchor $e<\len{w_l}$. Let $e$ be the last such anchor and
  $\mk tm$ the longest level-$t$ marker that ends at $e$. Then
  $\mk{t+1}l=\wsuf{w_l}{e-\len{\mk tm}}$.
\end{itemize}
The \emph{level-$0$ anchors} of a text are all its cuts $e\ge1$. The
level-$t$ anchors cut each nonempty word into its \emph{level-$t$ blocks},
the factors between consecutive anchors, the first block beginning at the
cut $0$. The \emph{top level} is the largest $t$ with $\Act_t\ne\varnothing$,
and $0$ if all words are empty.
\end{definition}

%% markers -- the marker hierarchy of the running example.
\begin{figure}[t]
\centering
\begin{tikzpicture}[font=\scriptsize, x=0.6cm, y=1cm]
\def\h{0.36}
%% #1 x-offset, #2 list of letters
\newcommand{\letters}[2]{%
  \foreach \c [count=\i from 0] in {#2} {
    \draw (#1+\i,0) rectangle (#1+\i+1,\h);
    \node at (#1+\i+0.5,0.5*\h) {$\c$};
  }}
%% a block row: #1 y, #2 x-offset, #3 list of cuts a/b
\newcommand{\blockrow}[3]{%
  \foreach \a/\b in {#3} {
    \draw[line width=0.9pt] (#2+\a+0.08,#1) -- (#2+\b-0.08,#1);
    \draw (#2+\a+0.08,#1-0.07) -- (#2+\a+0.08,#1+0.07);
    \draw (#2+\b-0.08,#1-0.07) -- (#2+\b-0.08,#1+0.07);
  }}
%% a marker above the word: #1 height, #2 x-offset, #3 from, #4 to, #5 label
\newcommand{\marker}[5]{%
  \draw[line width=1.5pt, black!55] (#2+#3+0.06,#1) -- (#2+#4,#1);
  \node[left, inner sep=1.5pt] at (#2+#3+0.06,#1) {#5};}
%% w1
\letters{0}{0,0,1,0,0,1,0,0,1,1}
\node[left] at (0,0.5*\h) {$w_1$};
\foreach \i in {0,...,10} { \node[below, font=\tiny, inner sep=1pt] at (\i,-0.02) {$\i$}; }
\marker{0.62}{0}{9}{10}{$\mk11=1$}
\marker{0.92}{0}{8}{10}{$\mk21=11$}
\marker{1.22}{0}{5}{10}{$\mk31=10011$}
\blockrow{-0.6}{0}{0/1,1/2,2/3,3/4,4/5,5/6,6/7,7/8,8/9,9/10}
\blockrow{-1.05}{0}{0/4,4/7,7/10}
\blockrow{-1.65}{0}{0/10}
\node[left] at (-0.3,-0.6) {level-$1$ blocks};
\node[left] at (-0.3,-1.05) {level-$2$ blocks};
\node[left] at (-0.3,-1.65) {level-$3$ blocks};
\foreach \x/\m in {4/10,7/10,10/11} { \node[below, inner sep=1pt] at (\x-0.3,-1.12) {$\m$}; }
%% w2
\letters{13.5}{1,1,0}
\node[left] at (13.5,0.5*\h) {$w_2$};
\foreach \i in {0,...,3} { \node[below, font=\tiny, inner sep=1pt] at (13.5+\i,-0.02) {$\i$}; }
\marker{0.62}{13.5}{2}{3}{$\mk12=0$}
\marker{0.92}{13.5}{1}{3}{$\mk22=10$}
\marker{1.22}{13.5}{0}{3}{$\mk32=110$}
\blockrow{-0.6}{13.5}{0/1,1/2,2/3}
\blockrow{-1.05}{13.5}{0/2,2/3}
\blockrow{-1.65}{13.5}{0/3}
\foreach \x/\m in {2/11,3/10} { \node[below, inner sep=1pt] at (13.5+\x-0.3,-1.12) {$\m$}; }
\end{tikzpicture}
\caption{The marker hierarchy of \cref{ex:hierarchy}. Above each word: its
markers of the levels $1$, $2$, $3$, nested suffixes of the word. Below: the
blocks of each level; a block ends where a marker of that level ends, in any
of the words. Under each level-$2$ block the level-$2$ marker that ends
there is shown: $10$ ends at the cuts $4$ and $7$ of $w_1$ and at the end of
$w_2$, and $11$ ends at the end of $w_1$ and at the cut $2$ of $w_2$.}
\label{fig:markers}
\end{figure}


So the blocks of level $0$ are single letters, the blocks of level $1$ end
with the last letter of some word, and the marker of the next level consists
of the last block of the word together with the marker in front of it. A
coordinate stops being active when its word is a single block.

\begin{example}\label{ex:hierarchy}
Let $W=(w_1,w_2)$ with $w_1=0010010011$ and $w_2=110$
(\cref{fig:markers}). The level-$1$
markers are the letters $1$ and $0$, so every cut is a level-$1$ anchor and
the level-$1$ blocks are single letters. In $w_1$ the last anchor before the
end is the cut $9$, where the marker $1$ ends, so $\mk21=11$; in $w_2$ it is
the cut $2$, and $\mk22=10$. The level-$2$ anchors are the cuts at which $11$
or $10$ ends: the cuts $4$, $7$ and $10$ of $w_1$, and the cuts $2$ and $3$
of $w_2$. They cut the words into the level-$2$ blocks
\[
  w_1=0010\cdot010\cdot011,\qquad w_2=11\cdot0 .
\]
At level $3$ the markers are $\mk31=10\,011$ (the level-$2$ marker $10$ that
ends at the cut $7$, followed by the last block) and $\mk32=11\,0=w_2$.
Neither occurs at another place, so the level-$3$ blocks are the whole words,
no coordinate is active at level $4$, and the top level is $3$.
\end{example}

The hierarchy has three structural properties. The first is the reason why
blocks behave well under comparisons.

\begin{lemma}[Markers]\label{lem:markers}
For every level $t\ge1$:
\begin{enumerate}
\item\label{it:mk-anti} \emph{Factor antichain.} If $l,m\in\Act_t$ and
  $\mk tl$ is a factor of $\mk tm$, then $\mk tl=\mk tm$. Consequently, if an
  occurrence of an active level-$t$ marker in a text lies inside an
  occurrence of another one, the two occurrences coincide.
\item\label{it:mk-two} For $l\in\Act_{t+1}$, the marker $\mk tl$ is a proper
  suffix of $\mk{t+1}l$, and the text $\mk{t+1}l$ has exactly two level-$t$
  anchors: the end of the level-$t$ marker it begins with, and its own end.
\item\label{it:mk-nest} Every level-$(t+1)$ anchor of a text is a level-$t$
  anchor of it.
\end{enumerate}
\end{lemma}

By \eqref{it:mk-nest} the blocks of level $t+1$ are concatenations of blocks
of level $t$. The freshness of $W$ enters through \cref{lem:distinct}: a
fixed text has at most $N-1$ nonempty rests of one word $w_l$ as prefixes.

\begin{lemma}[Depth and span]\label{lem:depth-span}
\begin{enumerate}
\item\label{it:depth} The top level is at most $k(N-1)$.
\item\label{it:span} In every text, an occurrence $\wsub uab$ of a
  level-$(t+1)$ marker contains at most $k(N-1)+2$ level-$t$ anchors $e$ of
  $u$ with $a<e\le b$.
\end{enumerate}
\end{lemma}

For \eqref{it:depth}, a marker of level $t$ has at least $t$ prefixes that
are suffixes of words, namely the markers nested at its left end, and these
give rests of words that are prefixes of one text. For \eqref{it:span}, a
level-$t$ marker that ends inside the occurrence lies inside it, which
happens at two anchors by \cref{lem:markers}, or it begins to the left of
it, and then the part inside is a rest of a word and a prefix of the marker.

\subsection{Parses and the Reverse Trie}\label{sec:fresh-trie}

Fix a level $t$ between $0$ and the top level. We call the blocks of level
$t+1$ \emph{upper blocks} and those of level $t$ \emph{lower blocks}; at the
top level the upper blocks are the nonempty words. An \emph{upper occurrence}
is an upper block together with its place, a word and two cuts. Its
\emph{parse} is the sequence of its lower blocks, a word over the alphabet of
lower blocks, and for $d\ge0$ its \emph{cut at depth $d$} is the cut that
lies $d$ lower blocks to the left of its end.

Equal texts are parsed alike only up to their left ends: a marker that
begins before an upper block can end inside it and create an additional
anchor there. Different occurrences of one text can therefore have different
parses, but by \cref{lem:distinct} only at boundedly many places, all to the
left of the first anchor that the text has on its own
(\cref{lem:parse-local}). What matters is how parses end. For two parses $Z$
and $Y$ let $\lcsuf(Z,Y)$ be the number of lower blocks in their longest
common suffix.

\begin{definition}[Reverse trie]\label{def:trie}
The \emph{nodes} on the path of a parse $Z$ are $0$, the length of $Z$, and
the numbers $\lcsuf(Z,Y)$ for the parses $Y$ of all upper occurrences. An
\emph{edge} is a parse $Z$ with two consecutive nodes $a<b$ on its path. An
upper occurrence \emph{traverses} the edge if its parse $Y$ satisfies
$\lcsuf(Y,Z)\ge b$.
\end{definition}

This is the trie of all parses read from the right; an edge is a stretch of
a parse along which no other parse branches off. All upper occurrences that
traverse an edge end with the same $b$ lower blocks, and whether an
occurrence traverses the edge is decided by a bounded part of the text at
its right end (\cref{lem:traverse}).

\begin{example}\label{ex:trie}
In \cref{ex:hierarchy} take $t=1$. The lower blocks are the letters, and the
upper blocks are the level-$2$ blocks $0010$, $010$, $011$ of $w_1$ and
$11$, $0$ of $w_2$; the parse of an upper block is the sequence of its
letters. For $Z=0010$ we have $\lcsuf(Z,010)=3$, $\lcsuf(Z,0)=1$ and
$\lcsuf(Z,011)=\lcsuf(Z,11)=0$. So the nodes on the path of $Z$ are
$0,1,3,4$, and its edges are $(Z,0,1)$, $(Z,1,3)$ and $(Z,3,4)$
(\cref{fig:trie}). The edge $(Z,1,3)$ is traversed by the upper blocks
$0010$ and $010$ of $w_1$, which end with the same three lower blocks
$0,1,0$; the edge $(Z,0,1)$ is traversed by these two and by the upper block
$0$ of $w_2$. The upper block $010$ lies between the cuts $4$ and $7$ of
$w_1$, and its cuts at the depths $0,1,2,3$ are $7,6,5,4$.
\end{example}

%% trie -- the reverse trie of the running example at level 1.
\begin{figure}[t]
\centering
\begin{tikzpicture}[font=\scriptsize, x=1.5cm, y=1cm,
    nd/.style={circle, draw, fill=white, inner sep=1.3pt},
    en/.style={circle, draw, fill=black, inner sep=1.3pt},
    pt/.style={circle, fill=black!40, inner sep=0.7pt}]
%% depth d is drawn at x = -d; root at the right
\node[nd] (r) at (0,0) {};
%% branch ending with 0: parses 0, 010, 0010
\node[en] (a1) at (-1,0.6) {};
\node[pt] (a2) at (-2,0.6) {};
\node[en] (a3) at (-3,0.6) {};
\node[en] (a4) at (-4,0.6) {};
\draw (r) -- node[above right, inner sep=1.5pt] {$0$} (a1);
\draw (a1) -- node[below] {$1$} (a2);
\draw (a2) -- node[below] {$0$} (a3);
\draw (a3) -- node[below] {$0$} (a4);
\node[above=2pt] at (a1) {$0$};
\node[above=2pt] at (a3) {$010$};
\node[left=2pt] at (a4) {$Z=0010$};
%% branch ending with 1: parses 11, 011
\node[pt] (b1) at (-1,-0.6) {};
\node[en] (b2) at (-2,-0.6) {};
\node[en] (b3) at (-3,-0.6) {};
\draw (r) -- node[below right, inner sep=1.5pt] {$1$} (b1);
\draw (b1) -- node[above] {$1$} (b2);
\draw (b2) -- node[above] {$0$} (b3);
\node[below=2pt] at (b2) {$11$};
\node[left=2pt] at (b3) {$011$};
%% edges of the path of Z
\foreach \a/\b/\l in {0/1/{(Z,0,1)},1/3/{(Z,1,3)},3/4/{(Z,3,4)}} {
  \draw[decorate, decoration={brace, amplitude=3pt}] (-\b+0.06,1.35) -- (-\a-0.06,1.35)
    node[midway, above=3pt] {$\l$};
}
%% depth axis
\draw[->, >={Stealth[length=1.8mm]}] (0,-1.65) -- (-4.5,-1.65) node[left] {depth};
\foreach \d in {0,...,4} { \draw (-\d,-1.58) -- (-\d,-1.72) node[below] {$\d$}; }
\end{tikzpicture}
\caption{The reverse trie of \cref{ex:trie}: the parses of the five upper
blocks, read from the right; the edges of the drawing are labeled with
lower blocks, here letters. A black node is the end of the parse written
next to it. The path of
$Z=0010$ has the nodes $0,1,3,4$; nothing branches off at depth $2$, so
$(Z,1,3)$ is one edge, which is traversed by the upper blocks $0010$ and
$010$.}
\label{fig:trie}
\end{figure}


\subsection{Block Deletions}\label{sec:fresh-deletion}

For a cut $c$ of a word $w_l$, let $\tau_r(l,c)$ be the text of the
level-$r$ marker that ends at $c$, if there is one; it is unique by
\cref{lem:markers}. Let $\omega(l,c)$ be the largest length of these texts
for $1\le r\le t$, and $0$ if there are none. The \emph{label} of the cut is
\[
  \lab(l,c)=\bigl(\tau_1(l,c),\dots,\tau_t(l,c),\;q_l(c-\omega(l,c))\bigr):
\]
the markers of all levels up to $t$ that end at the cut, and the state
before the longest of them. Two cuts of one word with equal labels have the
same $\omega$ letters to their left and the same states at all distances up
to $\omega$, in particular the same state. In \cref{ex:hierarchy} with
$t=2$, the markers $0$ of level $1$ and $10$ of level $2$ end at the cut $7$
of $w_1$, so $\omega(1,7)=2$ and $\lab(1,7)=(0,\,10,\,q_1(5))$; likewise
$\lab(1,4)=(0,\,10,\,q_1(2))$. The two labels are equal exactly when
$q_1(2)=q_1(5)$, and then also $q_1(3)=q_1(6)$ and $q_1(4)=q_1(7)$, because
the same letters $1,0$ are read from equal states.

Deletions will be restricted to the interior of an edge. The \emph{margin}
\[
  \Marg=3kN-k+5
\]
is the number of lower blocks that are kept free at both ends of the edge.
At the end $b$, where the parses branch, the margin ensures that the markers
ending at the cuts in question lie inside the text that all traversing
occurrences share, so that their labels can be compared
(\cref{lem:labels-local}). At the end $a$ it leaves enough common lower
blocks to the right of a deleted factor to recognize, from the text alone,
every upper block that traverses the edge (\cref{lem:traverse}).

\begin{definition}[Block deletion]\label{def:block-deletion}
A \emph{block deletion} at level $t$ consists of an edge $(Z,a,b)$ and two
depths $i<j$ with $a+\Marg\le i$ and $j\le b-\Marg$ such that, in every upper
occurrence $o$ that traverses the edge, the cuts of $o$ at the depths $j$ and
$i$ have equal labels. It deletes, from every traversing upper occurrence,
the factor between its cuts at the depths $j$ and $i$.
\end{definition}

\Cref{fig:hierarchy} illustrates the definition. The deleted factors are
equal words, since traversing occurrences end with the same $b$ lower
blocks, and at least one factor is deleted. The words of the example are
far too short for a block deletion, since an edge must have more than
$2\Marg$ lower blocks; in general, deletions exist as soon as some edge is
long (\cref{sec:fresh-bound}).

%% hierarchy -- a block deletion along an edge of the reverse trie, before and after.
\begin{figure}[t]
\centering
\begin{tikzpicture}[font=\scriptsize, >={Stealth[length=1.8mm]}]
\def\h{0.34}
%% draws an upper block ending at x=#1, y=#2, with a context-dependent left part of width #3
\newcommand{\upperblock}[4]{%
  % common suffix: blocks of depths 0..b-1 (b = 9 lower blocks of width .55)
  \foreach \d in {0,...,8} {
    \pgfmathsetmacro{\xr}{#1-0.55*\d}
    \pgfmathsetmacro{\xl}{#1-0.55*(\d+1)}
    \ifnum\d>2 \ifnum\d<6
      \fill[black!22] (\xl,#2) rectangle (\xr,#2+\h);
    \fi\fi
    \draw (\xl,#2) rectangle (\xr,#2+\h);
  }
  % the part above the edge, different in each occurrence
  \pgfmathsetmacro{\xe}{#1-0.55*9}
  \draw[pattern=north east lines, pattern color=black!35] (\xe-#3,#2) rectangle (\xe,#2+\h);
  \draw[line width=1.1pt] (#1,#2-0.08) -- (#1,#2+\h+0.08);
  \draw[line width=1.1pt] (\xe-#3,#2-0.08) -- (\xe-#3,#2+\h+0.08);
  \node[left] at (\xe-#3-0.05,#2+0.5*\h) {#4};
}
%% the same upper block after the deletion: 6 lower blocks of the edge remain
\newcommand{\upperblockafter}[4]{%
  \foreach \d in {0,...,5} {
    \pgfmathsetmacro{\xr}{#1-0.55*\d}
    \pgfmathsetmacro{\xl}{#1-0.55*(\d+1)}
    \draw (\xl,#2) rectangle (\xr,#2+\h);
  }
  \pgfmathsetmacro{\xe}{#1-0.55*6}
  \pgfmathsetmacro{\xj}{#1-0.55*3}
  \draw[pattern=north east lines, pattern color=black!35] (\xe-#3,#2) rectangle (\xe,#2+\h);
  \draw[line width=1.1pt] (#1,#2-0.08) -- (#1,#2+\h+0.08);
  \draw[line width=1.1pt] (\xe-#3,#2-0.08) -- (\xe-#3,#2+\h+0.08);
  \draw[line width=1.3pt, densely dashed] (\xj,#2-0.17) -- (\xj,#2+\h+0.17);
  \node[left] at (\xe-#3-0.05,#2+0.5*\h) {#4};
}
\upperblock{9.2}{2.0}{1.3}{$w_1$:}
\upperblock{9.2}{1.2}{0.5}{$w_1$:}
\upperblock{9.2}{0.4}{2.0}{$w_2$:}
%% depth scale
\draw[->] (9.2,2.75) -- (3.9,2.75) node[left] {depth};
\foreach \d/\l in {0/0,3/i,6/j,9/b} {
  \pgfmathsetmacro{\x}{9.2-0.55*\d}
  \draw (\x,2.68) -- (\x,2.82) node[above] {$\l$};
}
\draw[densely dotted] (9.2-0.55*3,2.68) -- (9.2-0.55*3,0.3);
\draw[densely dotted] (9.2-0.55*6,2.68) -- (9.2-0.55*6,0.3);
\draw[densely dotted] (9.2-0.55*9,2.68) -- (9.2-0.55*9,0.3);
\draw[decorate, decoration={brace, mirror, amplitude=4pt}] (9.2-0.55*9,0.28) -- (9.2,0.28)
  node[midway, below=4pt] {the last $b$ lower blocks: equal in all three occurrences (the edge)};
\node[align=left, anchor=west] at (9.6,1.35) {upper anchors\\ (thick)};
\draw[->] (9.6,1.0) -- (9.25,0.75);
\node[anchor=east] at (1.6,2.75) {before:};
%% after the deletion
\begin{scope}[yshift=-3.5cm]
\upperblockafter{9.2}{2.0}{1.3}{$u_1$:}
\upperblockafter{9.2}{1.2}{0.5}{$u_1$:}
\upperblockafter{9.2}{0.4}{2.0}{$u_2$:}
\draw[decorate, decoration={brace, mirror, amplitude=4pt}] (9.2-0.55*3,0.28) -- (9.2,0.28)
  node[midway, below=4pt] {the last $i$ lower blocks};
\node[above] at (9.2-0.55*3,2.0+\h+0.17) {junctions (dashed)};
\node[anchor=east] at (1.6,2.55) {after:};
\end{scope}
\end{tikzpicture}
\caption{A block deletion at level $t$, before and after. Three upper blocks
(blocks of level $t+1$), two in $w_1$ and one in $w_2$, are cut into lower
blocks (blocks of level $t$) by the anchors of level $t$. Their parses end
with the same $b$ lower blocks, so all three traverse the same edge of the
reverse trie; they differ above it (hatched). If the cuts at the depths $j$
and $i$ carry equal labels in every traversing upper block, the lower blocks
between these two depths (gray) are deleted in all of them at once. In the
result $(u_1,u_2)$ the edited upper blocks still end with equal texts: every
junction, the place where a factor was removed, is followed by the same $i$
lower blocks.}
\label{fig:hierarchy}
\end{figure}


\begin{theorem}[Block deletions]\label{thm:block-deletion}
If $W$ is a fresh tuple, the result of a block deletion is a strictly
shorter tuple that freshly dominates $W$. Hence a minimal fresh tuple admits
no block deletion, at any level.
\end{theorem}

The proof has three steps (\cref{app:fresh-deletion}). First, the equal
labels make every deleted factor a loop of the run, so endpoints are
preserved, no new state is visited and the result is fresh. Second, the
deletion \emph{preserves the anchors} of all levels up to $t+1$: a marker
occurrence of the edited words that runs over a deletion site has only a
short part to the left of the site, because of the antichain property and of
\cref{lem:markers}\eqref{it:mk-two}, and this part lies in the window that
the equal labels make identical on both sides of the deleted factor. So the
edited words are cut into the same blocks. Third, the channels: consider a
comparison of the edited tuple that runs over a deletion site of the host.
The $i$ lower blocks between the site and the end of its upper block survive
in the host, and they are copied into the target as long as the comparison
continues. If it continues to the end of the upper block, the target has an
upper block ending with the same text, which therefore traverses the same
edge, and so the target has a deletion site at the corresponding cut; the
deleted factor can be put back on both sides. If the comparison ends earlier,
the target is a short rest of its word, and the equal labels provide a
comparison of the original words with the same states. The place of a word
of the result where a factor was removed is called a \emph{junction}.

\subsection{The Bound}\label{sec:fresh-bound}

On an edge, the label of the cut at a depth $d$ is the same function of the
state at the top of the edge in every traversing occurrence, when $d$ is at
least $\Marg$ below the top: by \cref{lem:depth-span} the markers that end
at this cut lie in the common text of the edge. There are at most
$(k+1)^tN^N$ such functions. So an edge of more than
\[
  E=(k+1)^{kN+1}N^N+2\Marg+1
\]
lower blocks has two depths with the same function between its margins, and
admits a block deletion. In a minimal fresh tuple every edge is therefore
shorter. This bounds the number of lower blocks of an upper block by the
number of nodes on its path times $E$, hence by the number of parses. The
number of parses is at most $2^{k(N-1)}$ times the number of distinct texts
of upper blocks, and each text of a lower block is the text of an upper
block one level below. Starting from the top level, where there are at most
$k$ upper blocks, this gives a bound for every level.

\begin{theorem}[Fresh tuples, with the bound]\label{thm:fresh-bound}
Let $F=2^{k(N-1)}$, $\nu_0=k+1$ and $\nu_{d+1}=\nu_d+\nu_dF(\nu_dF+2)E$,
and let $X=(\nu_{kN}F+2)E$. Every minimal fresh tuple of a table with $k$
starts and $N$ states has total length at most
\[
  \FreshB(k,N)=k\,X^{kN+1}\;\le\;\expo_2\bigl((k+N+16)^4\bigr).
\]
\end{theorem}

Here $\nu_d$ bounds the number of distinct texts of upper blocks $d$
levels below the top level. In a minimal fresh tuple every upper block has
at most $X$ lower blocks at every level, there are at most $kN+1$ levels,
and the blocks of level $0$ are letters; this gives $k\,X^{kN+1}$.
